% GENERATED FILE. Edit canonical lessons, then run npm run generate:books.
% canonical-source-hash: fnv1a64:48a49d4e2228c87b
% canonical-lessons: TA-C480-mantai, TA-C480-tatai, TA-C480-puruvam, TA-C480-imai, TA-C480-cimittu

\chapter{Head and Face}
\label{ch:ta-head-and-face}
\hlchaptermodality{480}

\begin{chapteropening}
\textbf{By the end of this chapter:} \emph{I can talk about the skull, the jaw, eyebrows, eyelids and blinking.}

Complete the last lesson of chapter 480: I can talk about the skull, the jaw, eyebrows, eyelids and blinking.
\end{chapteropening}

\section[\texorpdfstring{\ta{மண்டை} (maṇṭai)}{மண்டை (maṇṭai)}]{\ta{மண்டை} (maṇṭai) — the skull; the head}

\label{lesson:TA-C480-mantai}



\begin{quote}
Before the new one: say the Tamil for a rib, then the Tamil for to throb.
\end{quote}

\subsection*{You'll want to know: \ta{மண்டை}}
\textbf{\ta{மண்டை}} — \emph{maṇṭai} — \textquotedblleft{}the skull; the head\textquotedblright{}.

\begin{quote}
\ta{வெயிலில்} \ta{நின்றதால்} \ta{எனக்கு} \ta{மண்டை} \ta{வலித்தது}.
\end{quote}

\begin{quote}
\emph{veyilil niṉṟatāl eṉakku maṇṭai valittatu.}
\end{quote}

\begin{quote}
Standing in the sun gave me a headache.
\end{quote}

Everyday and blunt for the head; \ta{தலை} is the neutral word.

\begin{cousinweb}
\textbf{\ta{மண்டை}} (\emph{maṇṭai}) is native Dravidian (DEDR 4682); the entry also gives the senses of a begging bowl and an earthen vessel.
\end{cousinweb}

\begin{grammarlens}[title={What you've built}]
One more word for describing what you lived and felt.
\end{grammarlens}

\subsection*{Your turn}
\begin{itemize}
\raggedright
  \item \emph{Say it:} \emph{maṇṭai}
  \item \emph{Say it:} \emph{maṇṭai}, once more
  \item \emph{Say it:} \emph{veyilil niṉṟatāl eṉakku maṇṭai valittatu.} — Standing in the sun gave me a headache.
  \item \emph{Recall:} say the Tamil for a rib, then the Tamil for to throb, then say \emph{maṇṭai} again
\end{itemize}

\begin{tcolorbox}[breakable,colback=teal!4,colframe=teal!35!black,title={Before you move on}]
What does \ta{மண்டை} mean? (\textquotedblleft{}The skull; the head\textquotedblright{}.) Say it once more.
\end{tcolorbox}
\section[\texorpdfstring{\ta{தாடை} (tāṭai)}{தாடை (tāṭai)}]{\ta{தாடை} (tāṭai) — the jaw}

\label{lesson:TA-C480-tatai}



\begin{quote}
Before the new one: say the Tamil for to throb, then the Tamil for the skull.
\end{quote}

\subsection*{You'll want to know: \ta{தாடை}}
\textbf{\ta{தாடை}} — \emph{tāṭai} — \textquotedblleft{}the jaw\textquotedblright{}.

\begin{quote}
\ta{பல்} \ta{வலியால்} \ta{என்} \ta{தாடை} \ta{வீங்கியது}.
\end{quote}

\begin{quote}
\emph{pal valiyāl eṉ tāṭai vīṅkiyatu.}
\end{quote}

\begin{quote}
My jaw swelled up because of a toothache.
\end{quote}

\ta{தாடை} and \ta{தாடி} \textquotedblleft{}beard\textquotedblright{} are different words.

\begin{cousinweb}
The origin of \textbf{\ta{தாடை}} (\emph{tāṭai}) is not settled; we could not find it in DEDR.
\end{cousinweb}

\begin{grammarlens}[title={What you've built}]
One more word for describing what you lived and felt.
\end{grammarlens}

\subsection*{Your turn}
\begin{itemize}
\raggedright
  \item \emph{Say it:} \emph{tāṭai}
  \item \emph{Say it:} \emph{tāṭai}, once more
  \item \emph{Say it:} \emph{pal valiyāl eṉ tāṭai vīṅkiyatu.} — My jaw swelled up because of a toothache.
  \item \emph{Recall:} say the Tamil for to throb, then the Tamil for the skull, then say \emph{tāṭai} again
\end{itemize}

\begin{tcolorbox}[breakable,colback=teal!4,colframe=teal!35!black,title={Before you move on}]
What does \ta{தாடை} mean? (\textquotedblleft{}The jaw\textquotedblright{}.) Say it once more.
\end{tcolorbox}
\section[\texorpdfstring{\ta{புருவம்} (puruvam)}{புருவம் (puruvam)}]{\ta{புருவம்} (puruvam) — an eyebrow}

\label{lesson:TA-C480-puruvam}



\begin{quote}
Before the new one: say the Tamil for the skull, then the Tamil for the jaw.
\end{quote}

\subsection*{You'll want to know: \ta{புருவம்}}
\textbf{\ta{புருவம்}} — \emph{puruvam} — \textquotedblleft{}an eyebrow\textquotedblright{}.

\begin{quote}
\ta{ஆச்சரியத்தில்} \ta{அவள்} \ta{புருவத்தை} \ta{உயர்த்தினாள்}.
\end{quote}

\begin{quote}
\emph{āccariyattil avaḷ puruvattai uyarttiṉāḷ.}
\end{quote}

\begin{quote}
She raised her eyebrows in surprise.
\end{quote}

The oblique -\ta{அத்து}-: \ta{புருவத்தை}.

\begin{cousinweb}
\textbf{\ta{புருவம்}} (\emph{puruvam}) is borrowed from Sanskrit \emph{bhrū}, \textquotedblleft{}eyebrow\textquotedblright{} (CDIAL 9688), with vowels breaking the cluster and the ending -\ta{அம்}.
\end{cousinweb}

\begin{grammarlens}[title={What you've built}]
One more word for describing what you lived and felt.
\end{grammarlens}

\subsection*{Your turn}
\begin{itemize}
\raggedright
  \item \emph{Say it:} \emph{puruvam}
  \item \emph{Say it:} \emph{puruvam}, once more
  \item \emph{Say it:} \emph{āccariyattil avaḷ puruvattai uyarttiṉāḷ.} — She raised her eyebrows in surprise.
  \item \emph{Recall:} say the Tamil for the skull, then the Tamil for the jaw, then say \emph{puruvam} again
\end{itemize}

\begin{tcolorbox}[breakable,colback=teal!4,colframe=teal!35!black,title={Before you move on}]
What does \ta{புருவம்} mean? (\textquotedblleft{}An eyebrow\textquotedblright{}.) Say it once more.
\end{tcolorbox}
\section[\texorpdfstring{\ta{இமை} (imai)}{இமை (imai)}]{\ta{இமை} (imai) — an eyelid}

\label{lesson:TA-C480-imai}



\begin{quote}
Before the new one: say the Tamil for the jaw, then the Tamil for an eyebrow.
\end{quote}

\subsection*{You'll want to know: \ta{இமை}}
\textbf{\ta{இமை}} — \emph{imai} — \textquotedblleft{}an eyelid\textquotedblright{}.

\begin{quote}
\ta{குழந்தையின்} \ta{இமைகள்} \ta{மெதுவாக} \ta{மூடின}.
\end{quote}

\begin{quote}
\emph{kuḻantaiyiṉ imaikaḷ metuvāka mūṭiṉa.}
\end{quote}

\begin{quote}
The baby’s eyelids slowly closed.
\end{quote}

Plural: \ta{இமைகள்}, \emph{imaikaḷ}.

The same word is a verb, to wink.

\begin{cousinweb}
\textbf{\ta{இமை}} (\emph{imai}) is native Dravidian (DEDR 2545), both the eyelid and the verb \textquotedblleft{}to wink\textquotedblright{}.
\end{cousinweb}

\begin{grammarlens}[title={What you've built}]
One more word for describing what you lived and felt.
\end{grammarlens}

\subsection*{Your turn}
\begin{itemize}
\raggedright
  \item \emph{Say it:} \emph{imai}
  \item \emph{Say it:} \emph{imai}, once more
  \item \emph{Say it:} \emph{kuḻantaiyiṉ imaikaḷ metuvāka mūṭiṉa.} — The baby’s eyelids slowly closed.
  \item \emph{Recall:} say the Tamil for the jaw, then the Tamil for an eyebrow, then say \emph{imai} again
\end{itemize}

\begin{tcolorbox}[breakable,colback=teal!4,colframe=teal!35!black,title={Before you move on}]
What does \ta{இமை} mean? (\textquotedblleft{}An eyelid\textquotedblright{}.) Say it once more.
\end{tcolorbox}
\section[\texorpdfstring{\ta{சிமிட்டு} (cimiṭṭu)}{சிமிட்டு (cimiṭṭu)}]{\ta{சிமிட்டு} (cimiṭṭu) — to blink; to wink}

\label{lesson:TA-C480-cimittu}



\begin{quote}
Before the new one: say the Tamil for an eyebrow, then the Tamil for an eyelid.
\end{quote}

\subsection*{You'll want to know: \ta{சிமிட்டு}}
\textbf{\ta{சிமிட்டு}} — \emph{cimiṭṭu} — \textquotedblleft{}to blink; to wink\textquotedblright{}.

\begin{quote}
\ta{அவன்} \ta{என்னைப்} \ta{பார்த்துக்} \ta{கண்} \ta{சிமிட்டினான்}.
\end{quote}

\begin{quote}
\emph{avaṉ eṉṉaip pārttuk kaṇ cimiṭṭiṉāṉ.}
\end{quote}

\begin{quote}
He looked at me and winked.
\end{quote}

Past (I …): \ta{சிமிட்டினேன்}, \emph{cimiṭṭiṉēṉ}.

\ta{கண்} \ta{சிமிட்டு}, to wink. \ta{இமை}, the eyelid, is in the same family.

\begin{cousinweb}
\textbf{\ta{சிமிட்டு}} (\emph{cimiṭṭu}) is native Dravidian (DEDR 2545), listed in the same entry as \ta{இமை} \textquotedblleft{}eyelid\textquotedblright{}.
\end{cousinweb}

\begin{grammarlens}[title={What you've built}]
One more word for describing what you lived and felt.
\end{grammarlens}

\subsection*{Your turn}
\begin{itemize}
\raggedright
  \item \emph{Say it:} \emph{cimiṭṭu}
  \item \emph{Say it:} \emph{cimiṭṭu}, once more
  \item \emph{Say it:} \emph{avaṉ eṉṉaip pārttuk kaṇ cimiṭṭiṉāṉ.} — He looked at me and winked.
  \item \emph{Recall:} say the Tamil for an eyebrow, then the Tamil for an eyelid, then say \emph{cimiṭṭu} again
\end{itemize}

\begin{tcolorbox}[breakable,colback=teal!4,colframe=teal!35!black,title={Before you move on}]
What does \ta{சிமிட்டு} mean? (\textquotedblleft{}To blink; to wink\textquotedblright{}.) Say it once more.
\end{tcolorbox}
